\clearpage\thispagestyle{empty}
\pdfbookmark[0]{Quellen zur Besprechung von Regners Aufsatz zur Sonnenparallaxe}{regner1810-sources}
\begin{center}\textbf{Quellen zur Besprechung von Regners Aufsatz zur Sonnenparallaxe}\end{center}
\noindent Diese KI-übersetzte Fassung in modernem Deutsch folgt Carl Friedrich Gauß, \href{https://gdz.sub.uni-goettingen.de/id/PPN236007467}{\emph{Werke}, Band~VI} (Göttingen, 1874), S.~523–525. Die Besprechung erschien zuerst ohne Verfassernamen unter Göttingen in den \emph{Göttingischen gelehrten Anzeigen}, Stück~17, 29.~Januar 1810, \href{https://images.p.sub.uni-goettingen.de/iiif/image/gdz:PPN319721507_1810:00000193/full/1650,/0/default.jpg}{S.~161}–\href{https://images.p.sub.uni-goettingen.de/iiif/image/gdz:PPN319721507_1810:00000198/full/1650,/0/default.jpg}{166}. Das \href{https://www.astro.uu.se/history/profflist_eng.html}{Verzeichnis der Sternwarte Uppsala} führt den Astronomen als Lars Regnér. In der Besprechung steht sein Name als Regner.

\noindent Der erste lateinische Titel bedeutet \emph{Ergänzung zur Geschichte der Entdeckung der Sonnenparallaxe, der hochangesehenen Königlichen Gesellschaft der Wissenschaften zu Göttingen zur Beurteilung vorgelegt}. Der zweite bedeutet \emph{Bemerkungen zum Bericht des Herrn Schubert an die Kaiserliche Akademie der Wissenschaften in Petersburg über eine Schrift zur Sonnenparallaxe}. Die lateinischen Titel folgen der Werkausgabe, einschließlich der hochgestellten Buchstaben nach \emph{D} im zweiten Titel. Im Erstdruck steht dieser Titel auf \href{https://images.p.sub.uni-goettingen.de/iiif/image/gdz:PPN319721507_1810:00000196/full/1650,/0/default.jpg}{S.~164}. Übersetzung und Satz mit KI erstellt.
\clearpage
